% GENERATED FILE. Edit canonical lessons, then run npm run generate:books.
% canonical-source-hash: fnv1a64:cfb3ad9aba0a1a55
% canonical-lessons: GU-C307-ugarvu, GU-C307-rni, GU-C307-udu, GU-C307-hemkhem, GU-C307-has

\chapter{You Saved Her Life}
\label{ch:gu-you-saved-her-life}
\hlchaptermodality{307}

\begin{chapteropening}
\textbf{By the end of this chapter:} \emph{I can thank someone who saved a life.}

Complete the last lesson of chapter 307: I can thank someone who saved a life.
\end{chapteropening}

\section[\texorpdfstring{\gu{ઉગારવું} (ugārvũ)}{ઉગારવું (ugārvũ)}]{\gu{ઉગારવું} (ugārvũ) — to rescue}

\label{lesson:GU-C307-ugarvu}



\begin{quote}
Before the new one: say the Gujarati for grain, then the Gujarati for to throng.
\end{quote}

\subsection*{You'll want to know: \gu{ઉગારવું}}
\textbf{\gu{ઉગારવું}} — \emph{ugārvũ} — \textquotedblleft{}to rescue\textquotedblright{}.

\begin{quote}
\gu{અરુણે} \gu{નદીમાંથી} \gu{એક} \gu{છોકરીને} \gu{ઉગારી}.
\end{quote}

\begin{quote}
\emph{Aruṇe nadīmā̃thī ek chokrīne ugārī.}
\end{quote}

\begin{quote}
Arun rescued a girl from the river.
\end{quote}

Past, with \gu{તેણે} (he or she …): \gu{તેણે} \gu{બાળકને} \gu{ઉગાર્યું}, \emph{teṇe bāḷakne ugāryũ}.

\begin{cousinweb}
\textbf{\gu{ઉગારવું}} (\emph{ugārvũ}) is the causative of \emph{ugarvũ}, to escape, which CDIAL lists under Sanskrit \emph{udgirati}, \textquotedblleft{}spits out, gives out\textquotedblright{} (CDIAL 1960): to make someone escape a danger. The person saved takes \textbf{\gu{ને}}: \textbf{\gu{છોકરીને} \gu{ઉગારી}}.
\end{cousinweb}

\begin{grammarlens}[title={What you've built}]
One more word for telling what happened.
\end{grammarlens}

\subsection*{Your turn}
\begin{itemize}
\raggedright
  \item \emph{Say it:} \emph{ugārvũ}
  \item \emph{Say it:} \emph{ugārvũ}, once more
  \item \emph{Say it:} \emph{Aruṇe nadīmā̃thī ek chokrīne ugārī.} — Arun rescued a girl from the river.
  \item \emph{Recall:} say the Gujarati for grain, then the Gujarati for to throng, then say \emph{ugārvũ} again
\end{itemize}

\begin{tcolorbox}[breakable,colback=teal!4,colframe=teal!35!black,title={Before you move on}]
What does \gu{ઉગારવું} mean? (\textquotedblleft{}To rescue\textquotedblright{}.) Say it once more.
\end{tcolorbox}
\section[\texorpdfstring{\gu{ઋણી} (ṛṇī)}{ઋણી (ṛṇī)}]{\gu{ઋણી} (ṛṇī) — indebted; grateful}

\label{lesson:GU-C307-rni}



\begin{quote}
Before the new one: say the Gujarati for to throng, then the Gujarati for to rescue.
\end{quote}

\subsection*{You'll want to know: \gu{ઋણી}}
\textbf{\gu{ઋણી}} — \emph{ṛṇī} — \textquotedblleft{}indebted; grateful\textquotedblright{}.

\begin{quote}
\gu{છોકરીની} \gu{માતાએ} \gu{અરુણને} \gu{કહ્યું}, \gu{હું} \gu{તમારી} \gu{ઋણી} \gu{છું}.
\end{quote}

\begin{quote}
\emph{chokrīnī mātāe Aruṇne kahyũ, hũ tamārī ṛṇī chũ.}
\end{quote}

\begin{quote}
The girl's mother told Arun that she was in his debt.
\end{quote}

\begin{cousinweb}
\textbf{\gu{ઋણી}} (\emph{ṛṇī}) is a tatsama, Sanskrit \emph{ṛṇin}, \textquotedblleft{}one who is in debt\textquotedblright{} (Monier-Williams), built on \emph{ṛṇa}, a debt. It agrees with the speaker: a man says \textbf{\gu{હું} \gu{તમારો} \gu{ઋણી} \gu{છું}}, a woman \textbf{\gu{તમારી} \gu{ઋણી}}.
\end{cousinweb}

\begin{grammarlens}[title={What you've built}]
One more word for telling what happened.
\end{grammarlens}

\subsection*{Your turn}
\begin{itemize}
\raggedright
  \item \emph{Say it:} \emph{ṛṇī}
  \item \emph{Say it:} \emph{ṛṇī}, once more
  \item \emph{Say it:} \emph{chokrīnī mātāe Aruṇne kahyũ, hũ tamārī ṛṇī chũ.} — The girl's mother told Arun that she was in his debt.
  \item \emph{Recall:} say the Gujarati for to throng, then the Gujarati for to rescue, then say \emph{ṛṇī} again
\end{itemize}

\begin{tcolorbox}[breakable,colback=teal!4,colframe=teal!35!black,title={Before you move on}]
What does \gu{ઋણી} mean? (\textquotedblleft{}Indebted; grateful\textquotedblright{}.) Say it once more.
\end{tcolorbox}
\section[\texorpdfstring{\gu{ઊંડું} (ū̃ḍũ)}{ઊંડું (ū̃ḍũ)}]{\gu{ઊંડું} (ū̃ḍũ) — deep}

\label{lesson:GU-C307-udu}



\begin{quote}
Before the new one: say the Gujarati for to rescue, then the Gujarati for indebted.
\end{quote}

\subsection*{You'll want to know: \gu{ઊંડું}}
\textbf{\gu{ઊંડું}} — \emph{ū̃ḍũ} — \textquotedblleft{}deep\textquotedblright{}.

\begin{quote}
\gu{કૂવો} \gu{ખૂબ} \gu{ઊંડો} \gu{હતો}.
\end{quote}

\begin{quote}
\emph{kūvo khūb ū̃ḍo hato.}
\end{quote}

\begin{quote}
The well was very deep.
\end{quote}

\begin{cousinweb}
\textbf{\gu{ઊંડું}} (\emph{ū̃ḍũ}) comes from \emph{*ūṇḍa}, \textquotedblleft{}deep\textquotedblright{}, a form CDIAL reconstructs, and CDIAL lists the Gujarati word (CDIAL 2402). It agrees like \textbf{\gu{મોટું}}: \textbf{\gu{ઊંડો} \gu{કૂવો}}, \textbf{\gu{ઊંડી} \gu{નદી}}, \textbf{\gu{ઊંડું} \gu{પાણી}}.
\end{cousinweb}

\begin{grammarlens}[title={What you've built}]
One more word for telling what happened.
\end{grammarlens}

\subsection*{Your turn}
\begin{itemize}
\raggedright
  \item \emph{Say it:} \emph{ū̃ḍũ}
  \item \emph{Say it:} \emph{ū̃ḍũ}, once more
  \item \emph{Say it:} \emph{kūvo khūb ū̃ḍo hato.} — The well was very deep.
  \item \emph{Recall:} say the Gujarati for to rescue, then the Gujarati for indebted, then say \emph{ū̃ḍũ} again
\end{itemize}

\begin{tcolorbox}[breakable,colback=teal!4,colframe=teal!35!black,title={Before you move on}]
What does \gu{ઊંડું} mean? (\textquotedblleft{}Deep\textquotedblright{}.) Say it once more.
\end{tcolorbox}
\section[\texorpdfstring{\gu{હેમખેમ} (hemkhem)}{હેમખેમ (hemkhem)}]{\gu{હેમખેમ} (hemkhem) — safe and sound}

\label{lesson:GU-C307-hemkhem}



\begin{quote}
Before the new one: say the Gujarati for indebted, then the Gujarati for deep.
\end{quote}

\subsection*{You'll want to know: \gu{હેમખેમ}}
\textbf{\gu{હેમખેમ}} — \emph{hemkhem} — \textquotedblleft{}safe and sound\textquotedblright{}.

\begin{quote}
\gu{બાળક} \gu{હેમખેમ} \gu{ઘરે} \gu{પહોંચ્યું}.
\end{quote}

\begin{quote}
\emph{bāḷak hemkhem ghare pahõcyũ.}
\end{quote}

\begin{quote}
The child got home safe and sound.
\end{quote}

\begin{cousinweb}
\textbf{\gu{હેમખેમ}} (\emph{hemkhem}) is an echo pair built on \emph{khem}, \textquotedblleft{}wellbeing\textquotedblright{}, which CDIAL lists for Gujarati under Sanskrit \emph{kṣ\'{\={e}}ma}, \textquotedblleft{}safety, ease\textquotedblright{} (CDIAL 3745); the rhyming \emph{hem} in front adds weight. It is said of people who come through danger: \textbf{\gu{હેમખેમ} \gu{પહોંચ્યું}}.
\end{cousinweb}

\begin{grammarlens}[title={What you've built}]
One more word for telling what happened.
\end{grammarlens}

\subsection*{Your turn}
\begin{itemize}
\raggedright
  \item \emph{Say it:} \emph{hemkhem}
  \item \emph{Say it:} \emph{hemkhem}, once more
  \item \emph{Say it:} \emph{bāḷak hemkhem ghare pahõcyũ.} — The child got home safe and sound.
  \item \emph{Recall:} say the Gujarati for indebted, then the Gujarati for deep, then say \emph{hemkhem} again
\end{itemize}

\begin{tcolorbox}[breakable,colback=teal!4,colframe=teal!35!black,title={Before you move on}]
What does \gu{હેમખેમ} mean? (\textquotedblleft{}Safe and sound\textquotedblright{}.) Say it once more.
\end{tcolorbox}
\section[\texorpdfstring{\gu{હાશ} (hāś)}{હાશ (hāś)}]{\gu{હાશ} (hāś) — phew (in relief)}

\label{lesson:GU-C307-has}



\begin{quote}
Before the new one: say the Gujarati for deep, then the Gujarati for safe and sound.
\end{quote}

\subsection*{You'll want to know: \gu{હાશ}}
\textbf{\gu{હાશ}} — \emph{hāś} — \textquotedblleft{}phew (in relief)\textquotedblright{}.

\begin{quote}
\gu{બાળક} \gu{બહાર} \gu{આવ્યું} \gu{ત્યારે} \gu{બધાએ} \gu{કહ્યું}, \gu{હાશ}!
\end{quote}

\begin{quote}
\emph{bāḷak bahār āvyũ tyāre badhāe kahyũ, hāś!}
\end{quote}

\begin{quote}
When the child came out, everyone breathed again: phew!
\end{quote}

\begin{cousinweb}
\textbf{\gu{હાશ}} (\emph{hāś}) No dictionary we use records an older form, so its history is not settled. It is the sound of breath let out after fear: \textbf{\gu{હાશ}!}, thank goodness.
\end{cousinweb}

\begin{grammarlens}[title={What you've built}]
One more word for telling what happened.
\end{grammarlens}

\subsection*{Your turn}
\begin{itemize}
\raggedright
  \item \emph{Say it:} \emph{hāś}
  \item \emph{Say it:} \emph{hāś}, once more
  \item \emph{Say it:} \emph{bāḷak bahār āvyũ tyāre badhāe kahyũ, hāś!} — When the child came out, everyone breathed again: phew!
  \item \emph{Recall:} say the Gujarati for deep, then the Gujarati for safe and sound, then say \emph{hāś} again
\end{itemize}

\begin{tcolorbox}[breakable,colback=teal!4,colframe=teal!35!black,title={Before you move on}]
What does \gu{હાશ} mean? (\textquotedblleft{}Phew (in relief)\textquotedblright{}.) Say it once more.
\end{tcolorbox}
